
49. Homomorphismes de $M_{0,3}$, les groupes $M_{g,\nu}$ généraux

I) Considérons un sous-groupe [unclear: fermé] $g$ de $M_{0,3}^\sim$, contenant $\widehat{\mathfrak{G}}_{0,3}^+$, donc image inverse d'un sous-groupe fermé de $\widehat{\Gamma}_{0,3}$.

[MARGIN: Oui on a un [unclear] Tout l'est... section [unclear] un ... le reste, [unclear] a un ... par un sous-groupe d'indice fini de [unclear] débauches p. 553]

On supposera que celui-ci est "d'épaisseur" grosse - disons qu'il contienne un sous-groupe d'indice fini de $\Gamma_{\mathbb{Q}} \subset \widehat{\Gamma}_{0,3}$. Pour le moment, on supposera pour simplifier que

$$g \subset K_{\varepsilon_3} . \quad \text{En résumé}$$

(1) $\widehat{\mathfrak{G}}_{0,3}^+ \subset g \subset M_{0,3}^\sim$, $u \in g \Rightarrow \mu(u) \equiv 1 \pmod 3$

Cette dernière hypothèse signifie donc que pour tout $u \in g$, $\alpha(u) \in \widehat{\mathfrak{G}}_{0,3}^+$, de sorte qu'on a

(2) $g \underset{\beta_0}{\overset{\alpha_0}{\rightrightarrows}} \widehat{\mathfrak{G}}_{0,3}^+$

caractérisés par

(2') $\begin{cases} u(\rho) = \text{int}(\alpha_0)\rho \\ u(\sigma) = \text{int}(\beta_0)\sigma \end{cases} \quad \begin{cases} \alpha_0 = \alpha(u) \\ \beta_0 = \beta(u) \end{cases}$

On considère

(3) $\begin{aligned} g_{(0)} &= g \cap M_{0,3}^\sim(0) \\ g_{[0]} &= g \cap M_{0,3}^\sim[0] = \text{Norm}_g(\mathbb{L}_0^\sim) \end{aligned}$

de sorte que

(4) $g = g_{(0)} \cdot \widehat{\mathfrak{G}}_{0,3}^+$ (produit semi-direct)

(5) $g_{[0]} = g_{(0)} \cdot \mathbb{L}_0^\sim$

(6) $g_{[0]} \cap \widehat{\mathfrak{G}}_{0,3}^+ = \mathbb{L}_0^\sim$

Où on a posé

(7) $\mathbb{L}_0^\sim =$ sous-groupe fermé de $\widehat{\mathfrak{G}}_{0,3}^+$ engendré par $\varepsilon_0$ i.e. [unclear] de $\mathbb{L}_1^\sim , \mathbb{L}_\infty^\sim$.

\newpage

%% ===== Page 496 =====

Considérons un [MARGIN: profini] groupe $M$ , et un hom. (continu) surjectif

(8) $$ \varphi : g \longrightarrow M .$$

On pose
(9) $$ \begin{cases} \varphi( \hat{\mathcal{G}}_{0,3}^+ ) = \mathcal{C} & \text{ss-groupe distingué de } M \\ \varphi( g_{(0)} ) = M_{(0)} \\ \varphi( g_{[0]} ) = M_{[0]} \\ \varphi( \tilde{L_0} ) = L_0 \end{cases} $$

Donc on aura les sous-groupes fermés de $M$, satisfaisant

(10) $$ M_{(0)} \subset M_{[0]} \subset Norm_M(L_0) $$
(11) $$ L_0 \subset \mathcal{C} \cap M_{[0]} $$

Par abus de langage, on désignera encore par $\rho, \sigma, \varepsilon_0, \varepsilon_1$ ... les images de ces éléments dans $\mathcal{C}$, qui y satisfont donc aux relations connues dans $\hat{\mathcal{G}}_{0,3}^+$ , en particulier

(12) $$ \sigma^2 = \rho^3 = 1 , \quad \sigma \rho = \varepsilon_0 , \quad \rho \sigma = \varepsilon_1 , \quad \varepsilon_1 = \sigma(\varepsilon_0) = \rho(\varepsilon_0) .$$

Les restrictions des opérations $\alpha_0, \beta_0$ de (2, 2') à $\mathcal{M}_{0,3}[0]$ seront plutôt notées $\alpha, \beta$ , où $\alpha, \beta$ sont en fait des opérations

(13) $$ Sg \overset{\alpha}{\underset{\beta}{\rightrightarrows}} \hat{\mathcal{G}}_{0,3}^+ $$
$$ Sg \overset{déf}{=} \text{ ext. par } \tilde{L_0} \text{ s-g de } S\mathcal{M}_{0,3} $$
$$ = \{ (u,f) \in g \times \hat{\mathcal{G}}_{0,3}^+ \mid f u \in g_{[0]} \} $$

dont on regarde les restrictions au sous-groupe ($\simeq \mathcal{M}_{0,3}[0]$) de $Sg$ défini par la relation $f=1$.

Comme $\alpha, \beta : g_{[0]} \rightrightarrows \hat{\mathcal{G}}_{0,3}^+$ ... se factorisent par $\varphi | \hat{\mathcal{G}}_{0,3}^+$ , donnent des opérations, encore notées $\alpha, \beta$ (i.e. $\alpha^M, \beta^M$ si on a une confusion possible)

(14) $$ g_{[0]} \overset{\alpha}{\underset{\beta}{\rightrightarrows}} \mathcal{C} $$

satisfaisant aux relations

(15) $$ \begin{cases} \alpha(p) = \rho \alpha(p) \rho \\ \beta(p) = \rho \beta(p) \sigma \end{cases} \quad \text{pour } u \in M_{[0]} \text{ (} \text{[unclear: image de]} \ \varphi(u_0)\text{),} $$
$$ \text{où } \alpha = \alpha(u_0), \beta = \beta(u_0). $$

\newpage

%% ===== Page 497 =====

A priori, $\alpha$ et $\beta$ dépendent des choix de $u_0$ dans ce procès
pas seulement de $u$ - mais
[MARGIN: (16)]
$$ \begin{cases} \alpha p(\alpha)^{-1} = [u, \rho] \\ \beta p(\beta)^{-1} = [u, \sigma] \end{cases} $$
n'en dépendent pas. Si on pose
[MARGIN: (17)]
$$ T_\rho = Cent_M(\rho) \qquad T_\sigma = Cent_M(\sigma) $$
on voit donc que, modulo $T_\rho$ resp. $T_\sigma$, $\alpha, \beta$ ne dépendent
que de $u$.

Hypothèse $\alpha (= \alpha(u_0))$, $\beta (= \beta(u_0))$ ne dépendent que de $u = \varphi(u_0)$ (pour $u_0 \in G[0]$)
On exprimera cette hypothèse de façon un peu
différente, en introduisant
[MARGIN: (18)]
$$ J = \ker(G[0] \longrightarrow M[0]) , $$
on veut donc que, pour
$u'_0 = u_0 \delta_0$ avec $u_0 \in G[0]$, $\delta_0 \in J$, on ait
$\alpha(u'_0) = \alpha(u_0)$, $\beta(u'_0) = \beta(u_0)$.
Soient $\alpha_0 = \alpha(u_0)$, $\beta_0 = \beta(u_0)$, $\alpha'_0 = \alpha(u'_0)$, $\beta'_0 = \beta(u'_0)$, $\bar{\alpha}_0 = \alpha(\delta_0)$, $\bar{\beta}_0 = \beta(\delta_0)$
dans $G$,
[unclear: d'après ce qu'on a vu], par les formules de cocycles de $G$ :
$$ \alpha'_0 = \delta_0(\alpha_0) \bar{\alpha}_0 , \quad \beta'_0 = \delta_0(\beta_0) \bar{\beta}_0 $$
d'où, en appliquant $\varphi$ et notant que
$$ \varphi(\delta_0(\alpha_0)) = \varphi(\delta_0) (\varphi(\alpha_0)) = \alpha \quad \text{et de même} = $$
$$ \varphi(\delta_0(\beta_0)) = \varphi(\delta_0) (\varphi(\beta_0)) = \beta $$
[MARGIN: (19)]
$$ \alpha' = \alpha \bar{\alpha} , \quad \beta' = \beta \bar{\beta} \quad \text{dans } \mathfrak{G} $$
où $\alpha, \beta, \alpha', \beta', \bar{\alpha}, \bar{\beta}$ sont les images par $\varphi$ des quantités
$\alpha_0$ etc. de $\widehat{\mathfrak{G}_{0,s}}^{+\wedge}$. Donc l'hypothèse
se [unclear: traduit alors] :
[MARGIN: (20)]
$$ \forall \delta_0 \in J = \ker(G[0] \to M[0]) , \quad \bar{\alpha}(\delta_0), \bar{\beta}(\delta_0) \in \ker p $$
Quand cette hypothèse est satisfaite, on a donc des
applications.

\newpage

%% ===== Page 498 =====

(21) 
$$ \mathcal{M}[\sigma_0] \overset{\alpha}{\underset{\beta}{\rightrightarrows}} \hat{G} $$
[unclear: satisfaisant] [crossed out]
précisant (14), et satisfaisant aux

(22)
$$ \begin{cases} u(\rho) = \alpha(u) \rho \\ u(\sigma) = \beta(u) \sigma \end{cases} \quad \text{pour } u \in \mathcal{M}[\sigma_0], \, \alpha = \alpha(u), \, \beta = \beta(u) , $$

qui généralisent (2').

D'ailleurs, posant

(23)
$$ \xi = \alpha \rho \alpha^{-1}, \quad \eta = \beta \sigma \beta^{-1} $$
les relations (22) s'écrivent aussi

(24)
$$ u(\rho) = \xi \rho, \quad u(\sigma) = \eta \sigma $$
où on a

(25)
$$ \xi = [u, \rho], \quad \eta = [u, \sigma] $$
et de plus

(26)
$$ \xi = \eta^\nu \quad \text{où} \quad \nu = \frac{\mu - 1}{2} \in \hat{\mathbb{Z}} \quad \text{(cf le multipli-} $$
[MARGIN: cateur de $u$).]

qui provient de la relation analogue dans $\hat{G}_{0,3}^+$.
(On suppose pour fixer les idées que
l'hom.

$$ g \xrightarrow{\chi^g} \hat{\mathbb{Z}}^\times \quad \text{("multiplicateur")} $$
se factorise par

(27)
$$ g \rightarrow \mathcal{M} \xrightarrow{\chi^\mathcal{M}} \hat{\mathbb{Z}}^\times , $$
de sorte qu'on a une notion de multiplicateur ($\mu \in \hat{\mathbb{Z}}^\times$)
pour les $u \in \mathcal{M}$.)

Considérons l'hom. naturel

(28)
$$ \mathcal{M}[\sigma_0] \rightarrow Aut(\hat{G}) $$
$$ u \mapsto int(u) | \hat{G} $$
on voit sur (24) qu'il est connu quand on connaît

\newpage

%% ===== Page 499 =====

[MARGIN: à part topologie de $F$ dont] les fonctions $\alpha, \beta : M[T_0] \to \tilde{\mathbb{G}}$, et par suite
en fait seulement $\xi, \eta : M[T_0] \to \tilde{\mathbb{G}}$.
D'ailleurs on pose

(29) $$M'[T_0] = Im(M[T_0] \to Aut(\tilde{\mathbb{G}}))$$

ces [unclear: mêmes] fonctions $\xi, \eta$ sont déjà définies sur
$M'[T_0]$ - alors qu'il n'est pas clair qu'il en soit
de même de $\alpha, \beta$, quand $M[T_0]$ n'opère pas fidèlement
sur $\tilde{\mathbb{G}}$. On a ainsi une application [unclear] de
certain sous-groupe $M'[T_0] \subset Norm_{Aut(\tilde{\mathbb{G}})}(L_0)$ dans
l'ensemble des solutions de l'équation en $\xi, \eta$
(26) qui s'écrit [unclear]

(30) $$\xi^{-1} \eta \in L_1$$

relation qui n'est que (26) (avec la forme précise de
l'exposant $\nu$) - [unclear] que
$\tilde{\mathbb{G}}_{0,s}^+ \to Sl(2, \hat{\mathbb{Z}})$
se factorise par $\tilde{\mathbb{G}}$ - ce qu'on
va supposer, pour nous simplifier la vie.

On va faire sur $\tilde{\mathbb{G}}$, muni de $\rho, \sigma$, l'hypo-
thèse [unclear] suivante

$$ \bar{X} = X_{\tilde{\mathbb{G}},\rho,\sigma} \subset \tilde{\mathbb{G}} \times \tilde{\mathbb{G}} $$
(31)
$$ \left\{ \begin{array}{l} (\xi, \eta) \text{ resp. conjugués de } \rho, \sigma \text{ dans } \tilde{\mathbb{G}} \\ \xi \eta^{\nu} \in L_1, \text{ i.e. } \exists l \in L_1 \text{ que } \xi = \eta^{-\nu} l \end{array} \right. $$

[unclear crossed out text]
Pour $u \in Aut(\tilde{\mathbb{G}})$ [unclear]
[unclear crossed out text] fixant les classes de conjugaison de $\rho$
et celles de $\sigma$, i.e. telles que
$\xi = [u, \rho]$, $\eta = [u, \sigma]$
satisfont aux deux premières conditions de (31), et

\newpage

%% ===== Page 500 =====

et soit $u \in Aut (\tilde{G})$ un automorphisme de $\tilde{G}$,
et posons

(34)
$$ \begin{cases} \xi = [u,\rho] = u \rho u^{-1} \rho^{-1} = u(\rho) \rho^{-1} & \text{i.e. } u(\rho) = \xi \rho \\ \eta = [u,\sigma] = u(\sigma) \sigma^{-1} & \text{i.e. } u(\sigma) = \eta \sigma \end{cases} $$
[MARGIN: cf (24) (25)]

alors dire que $\xi \rho$ $(=u(\rho))$ est conjugué à $\rho$ signifie que
$u$ fixe la cl. de c-j. de $\rho$, dire que $\eta \sigma$ $(=u(\sigma))$
conjugué de $\sigma$ signifie que $u$ fixe la classe de
c-j. de $\sigma$. Enfin la condition $\eta^{-1} \xi \in L_1$ s'écrit
encore, comme

[MARGIN: NB $\sigma^{-1} = \sigma$]

$$ \eta^{-1} \xi = (\sigma u(\sigma^{-1})) (u(\rho) \rho^{-1}) = \sigma u(\sigma^{-1} \rho) \rho^{-1} = \sigma (u(\sigma^{-1} \rho) \underbrace{\rho^{-1} \sigma}_{\xi_0}) \underbrace{\sigma^{-1} \rho^{-1}}_{\xi_0^{-1}} $$
$$ = \sigma (u(\xi_0) \xi_0^{-1}) $$

dans la forme
$$ \sigma (u(\xi_0) \xi_0^{-1}) \in L_1 \quad \text{i.e.} \quad u(\xi_0) \xi_0^{-1} \in L_0 $$
$(L_0 = \sigma^{-1} L_1)$ ) — ou encore (puisque $\xi_0 \in L_0$ [unclear: engendre $L_0$])
$u(\xi_0) \in L_0$, soit (puisque $u$ [unclear: conserve] $L_0$)
$$ u(L_0) = L_0 $$
i.e. que $u$ normalise $L_0$.

[MARGIN: [unclear: Dans notre cas où le $\xi_0$ (qui figure dans $\sigma \rho = \xi_0 \in L_1$) engendre $L_0$, de sorte que la condition $u(\xi_0) \in L_0$ équivaut à $u(L_0) = L_0$ ... $\rho^2=1$ ...]]

Soit donc

(35)
$$ \tilde{B}'_0 = Aut(\tilde{G}) $$
$$ \| dfn $$
$$ \{u \in Aut(\tilde{G}) \mid u \text{ fixe classes de c-j. de } \rho \text{ et celle de } \sigma, u \text{ normalise } L_0 \} $$

de sorte qu'on a une application in-jective pourvu que
[MARGIN: pour $\rho,\sigma$ eng. $\tilde{G}$]

(36)
$$ \tilde{B}'_0 \hookrightarrow \Sigma $$
$$ u \longmapsto ([u,\rho] = u(\rho)\rho^{-1}, [u,\sigma] = u(\sigma)\sigma^{-1}) $$

\newpage

%% ===== Page 501 =====

553

Hypothèse fondamentale sur $\overline{\mathcal{G}}$, muni de $\rho, \sigma$, et d'un $\varepsilon_0 = \sigma^{-1}\rho$ ,
$\varepsilon_1 = \rho \sigma^{-1} = \sigma(\varepsilon_0) = \rho(\varepsilon_0)$ :

[MARGIN: (37)]
a) $\overline{\mathcal{G}}$ est engendré topologiquement par $\rho, \sigma$
b) L'application (36) est (injective par a) et bijective.

Cela signifie donc que [unclear: text crossed out] $\Sigma$ est
La condition b) signifie donc que pour tout
$(\alpha, \beta) \in \overline{\mathcal{G}} \times \overline{\mathcal{G}}$, satisfaisant une relation
$$\alpha \rho(\alpha)^{-1} = \beta \sigma(\beta)^{-1} \varepsilon_1^{-1} \quad \text{s'écrivant aussi,}$$
il existe $u \in Aut(\overline{\mathcal{G}})$ tel qu'on ait
$$[u, \rho] = \alpha \rho(\alpha)^{-1} \quad \text{i.e.} \quad u(\rho) = int(\alpha)(\rho) \quad \text{i.e.} \quad \text{[unclear: } u \rho u^{-1} = \alpha \rho \alpha^{-1}]$$
$$[u, \sigma] = \beta \sigma(\beta)^{-1} \quad \text{i.e.} \quad u(\sigma) = int(\beta)(\sigma) \quad \text{i.e.} \quad \text{[unclear: } u \sigma u^{-1} = \beta \sigma \beta^{-1}]$$
(on aura donc nécessairement $u \in \widetilde{\mathcal{B}}'_0$, et
$u$ sera unique par a) ).

Ainsi, sur l'ens. $R_{\rho, \sigma}$ des $(\xi, \eta)$ satisfaisant
a) b) c) , on aura une structure de groupe,
déduite par transport de structure de celle de $\widetilde{\mathcal{B}}'_0$
[unclear]

\newpage

%% ===== Page 502 =====

[MARGIN: Les groupes $\mathcal{U}_{\rho,\sigma}$ etc]

[MARGIN: Pour une justification des notations qui suivent, cf [unclear: plus loin] p. 564, remarque.]

